% Hva indeksen kan gi informasjon om framover

\begin{frame}{Hva indeksen kan fange opp framover}
\begin{itemize}\setlength{\itemsep}{8pt}
  \item \textbf{Løpende overvåking.} Oppdateres månedlig etter hvert som nye A-ordningsdata publiseres, med rundt tre måneders datalag.
  \item \textbf{Klart signalkriterium.} Vintage-øvelsen viser at hovedtallet svinger 1--2 prosentpoeng fra måned til måned. Et reelt skift vil vise seg som et vedvarende utslag utover usikkerhetsbåndet, ikke som en enkeltobservasjon.
  \item \textbf{Flere serier enn hovedtallet.} Egne serier per tiårs aldersgruppe og for utvalgte yrker, der eventuelle utslag ventes først.
  \item \textbf{Re-ankring ved kapabilitetssprang.} Referansemåneden kan flyttes når teknologien tar nye steg, slik vi gjør for agentisk KI (april 2025).
\end{itemize}
\end{frame}

% ---------------------------------------------------------------
\begin{frame}{Avgrensninger og mulige utvidelser}
\begin{columns}[T,onlytextwidth]
\begin{column}{0.48\textwidth}
\textbf{\textcolor{illInk}{Hva indeksen ikke svarer på}}

\vspace{0.5em}
\small
\begin{itemize}\setlength{\itemsep}{6pt}
  \item Et beskrivende vekstgap, ikke en årsakseffekt av KI.
  \item Privat sektor; offentlig sektor inngår ikke i hovedtallet.
  \item Observerer arbeidsmarkedsutfall, ikke bruk eller oppgaveendring, jf.\ de fire størrelsene.
  \item Lønn er foreløpig lite informativt: sesongmønstre dominerer, og lønnsdannelsen demper utslag.
\end{itemize}
\end{column}
\begin{column}{0.48\textwidth}
\textbf{\textcolor{illInk}{Utvidelser som vurderes}}

\vspace{0.5em}
\small
\begin{itemize}\setlength{\itemsep}{6pt}
  \item Utdanningsgrupper, blant annet nyutdannede.
  \item Offentlig sektor.
  \item Arbeidssøkerdata.
  \item Andre eksponeringsmål.
  \item Med mer, gjerne i dialog med brukere av indeksen.
\end{itemize}
\end{column}
\end{columns}
\end{frame}
